% TOP_PROBLEM: {"id": 9500002, "title": "Topology of planar Brownian trace", "category_id": 19, "category": {"id": 19, "name": "probability", "display_name": "Probability", "description": "Problems involving probability theory, stochastic processes, and random structures.", "slug": "probability", "order_index": 19, "created_at": "2026-07-31T15:26:25.671Z"}, "status": "open", "problem_number": "AMR-094-0002", "set_id": 13}
% FIRST_POSTED: 2026-10-01
% ENTRY_KIND: statement_only
% UnsolvedMath ID 9500002; source catalogue code AMR-094-0002
% !TeX program = lualatex
% Definition and sources only; this file is not an LLM solution attempt.
\documentclass[11pt,a4paper]{article}
\usepackage[margin=25mm]{geometry}
\usepackage{fontspec}
\setmainfont{lmroman10-regular.otf}[BoldFont=lmroman10-bold.otf,
 ItalicFont=lmroman10-italic.otf,BoldItalicFont=lmroman10-bolditalic.otf]
\usepackage{amsmath,amssymb,mathtools}
\usepackage{unicode-math}
\setmathfont{latinmodern-math.otf}
\AtBeginDocument{\let\setminus\smallsetminus}
\usepackage{xurl,hyperref}
\hypersetup{colorlinks=true,urlcolor=blue}
\setcounter{secnumdepth}{0}
\setlength{\parindent}{0pt}
\setlength{\parskip}{5pt}
\setlength{\emergencystretch}{3em}
\begin{document}
\section*{Topology of planar Brownian trace}\label{9500002}
{\small UnsolvedMath ID 9500002 · AMR-094-0002 · Probability · AMR Open Problem Lists}
\subsection{Definitions and mathematical statement}
Let $X$ be standard planar Brownian motion on $[0,1]$, and let $T=X[0,1]=\{X_t:0\le t\le1\}$ be its compact trace. A Jordan arc is the image of a continuous injective map from $[0,1]$ into $\mathbb R^2$; equivalently, it is a set homeomorphic to that interval.

Two questions concern the almost-sure topology of this random set. First, does almost every trace have the property that for every pair of distinct points $x,y\in\mathbb R^2\setminus T$, there is a Jordan arc $\Gamma$ containing both points for which $\Gamma\cap T$ is finite? The arc and the finite number of intersections may depend on both the trace and the chosen points; no common bound is requested.

Second, enumerate the connected components of the open complement as $A_1,A_2,\ldots$, including its unbounded component, and set
\[
K=T\setminus\bigcup_{j\ge1}\partial A_j.
\]
Is $K$ almost surely totally disconnected, meaning that every connected subset of $K$ contains at most one point?

Boundaries are taken in the ordinary Euclidean topology. The complement has at most countably many components, so the displayed union is unambiguous. Both assertions are pathwise statements: the first quantifies over all exterior point pairs on one probability-one event, rather than allowing a different exceptional null event for each deterministic pair. The second concerns trace points which lie on none of the complementary component boundaries.
\subsection{Short English statement}
Can every two points outside a planar Brownian trace be joined through it only finitely often, and is the trace’s remaining interior topological set totally disconnected?
\subsection{Sources}
\begin{itemize}
\item[S1] ULAM.AI and the UnsolvedMath Contributors, supplied problems.json, record 9500002 (AMR-094-0002), AMR Open Problem Lists. Catalogue identity and source transcription; CC BY 4.0.
\url{https://huggingface.co/datasets/ulamai/UnsolvedMath}.
\item[S2] Burdzy - My favorite open problems (2009), Problem 2 (parts i and ii). Original source indexed by AMR-094-0002.
\url{https://sites.math.washington.edu/~burdzy/open_mathjax.php}.
\end{itemize}
\end{document}
